% Compile: pdflatex → bibtex → pdflatex → pdflatex
\documentclass{article}
\usepackage{neurips_2026}

\usepackage[utf8]{inputenc}
\usepackage[T1]{fontenc}
\usepackage{hyperref}
\usepackage{url}
\usepackage{booktabs}
\usepackage{amsfonts}
\usepackage{amsmath}
\usepackage{amssymb}
\usepackage{nicefrac}
\usepackage{microtype}
\usepackage{xcolor}
\usepackage{graphicx}
\usepackage{subcaption}
\usepackage{algorithm}
\usepackage{algpseudocode}
\usepackage{multirow}
\usepackage{makecell}

% Convenience colour macros
\definecolor{improvegreen}{RGB}{0,128,0}
\definecolor{worsered}{RGB}{180,0,0}

\title{Cross-Cultural Moral Alignment of Large Language Models\\
via Stochastic Welfare Agents and\\
Model Predictive Path Integral Control}

\author{}

\begin{document}

\maketitle

% ============================================================
\begin{abstract}
% ============================================================
Large language models (LLMs) exhibit a culturally monolithic moral stance that
systematically misaligns with the diverse ethical preferences of populations
worldwide. We propose \textbf{SWA-MPPI}-\textbf{S}tochastic \textbf{W}elfare
\textbf{A}gent with \textbf{M}odel \textbf{P}redictive \textbf{P}ath
\textbf{I}ntegral control-a training-free, inference-time framework that steers
LLM moral decisions toward country-specific human preferences without modifying
model weights. For each target country, SWA-MPPI constructs a small ensemble of
culturally grounded \emph{persona agents} from World Values Survey (WVS) Wave~7
data, evaluates their decision-level logit consensus via a single batched forward
pass, and applies a Prospect-Theory-weighted MPPI optimisation step whenever
inter-agent disagreement exceeds an adaptively calibrated threshold~$\tau$. We
evaluate on the Multilingual Trolley Problem benchmark (MultiTP), which provides
country-specific Average Marginal Component Effects (AMCE) derived from over
40~million human judgments across 233 countries and six moral dimensions. Experiments
with two frontier 70B-class models-Meta-Llama-3.1-70B-Instruct and
Qwen2.5-72B-Instruct-spanning 15 countries and 12 native languages show that
SWA-MPPI reduces Jensen-Shannon Distance to human moral preferences by
\textbf{37.6\%} (Qwen) and \textbf{34.7\%} (Llama) on average relative to a
vanilla LLM baseline, while improving Pearson correlation with human AMCE vectors
by up to $+0.35$ points. Alignment gains reach 66\% for culturally distant
country--language pairs. The entire framework operates on frozen weights, requires
no labelled preference data, and introduces a mean per-scenario overhead of
approximately 730\,ms on a single NVIDIA H100 GPU.
\end{abstract}

% ============================================================
\section{Introduction}
\label{sec:intro}
% ============================================================

As large language models are deployed across increasingly diverse cultural contexts,
a fundamental tension has emerged between the uniformity of a single trained model
and the plurality of human moral values worldwide. Prior work has established that
LLMs trained predominantly on English-language corpora inherit the ethical
intuitions of Western, educated, industrialised, rich, and democratic (WEIRD)
populations~\citep{henrich2010weirdest}, yet these systems are increasingly
enlisted to assist with consequential decisions in societies whose moral norms
differ substantially~\citep{awad2018moral, jin2025multitp}.

The Moral Machine experiment~\citep{awad2018moral} and its multilingual successor,
the Multilingual Trolley Problem dataset (MultiTP)~\citep{jin2025multitp}, provide
a large-scale, cross-culturally validated benchmark for studying these differences.
Built on over 40~million human judgments from 233 countries, the dataset captures
country-specific Average Marginal Component Effects (AMCE) across six moral
dimensions-species preference, gender, age, physical fitness, social status, and
utilitarianism. Against this ground truth, \citet{jin2025multitp} evaluated 19
contemporary LLMs and found that only the Llama~3 family achieves global
misalignment scores below 0.6 on the $[0,\sqrt{6}]$ scale they propose, while
many models-including GPT-4o Mini-show extreme preference binarisation that
departs markedly from the graded distributions observed in human respondents.

Existing approaches to value alignment fall broadly into two families.
\emph{Fine-tuning} on preference data~\citep{bai2022constitutional,
ouyang2022instructgpt} can improve cultural fit but requires large per-country
labelled datasets, modifies model weights irreversibly, and scales poorly to the
130-plus countries represented in MultiTP. \emph{Prompting-based steering} is
lightweight but typically relies on ad-hoc system instructions with no grounding
in empirical population-level values, and its effects on moral reasoning are
inconsistent across languages~\citep{jin2025multitp}.

We propose \textbf{SWA-MPPI}, a principled, training-free inference-time framework
that closes this gap. The method is motivated by three observations. First, cultural
moral preferences are not monolithic: within any society, age cohorts and ideological
subgroups exhibit systematically different moral weights~\citep{awad2018moral,
kleiman2017learning}. Second, the World Values Survey provides rich, empirically
calibrated data on exactly these sub-group differences in 80+ countries. Third,
when persona agents derived from this data disagree on a decision, that disagreement
is itself a signal that the vanilla model's output may not reflect the target
population's balance of views-a signal that can be leveraged through stochastic
optimal control.

SWA-MPPI proceeds as follows. For a target country, it builds four system-prompt
personas from WVS Wave~7 data (three age-cohort profiles plus a utilitarian agent).
It evaluates all agents in a single batched forward pass and extracts decision
logits for the tokens \texttt{LEFT} and \texttt{RIGHT}. If the variance of
agent decision gaps exceeds an adaptively calibrated threshold~$\tau$, it samples
$K = 128$ perturbations around the agent consensus signal and scores each perturbation
with a Prospect-Theory value function~\citep{kahneman1979prospect}, trading off
private and social utility with cooperation weight $\lambda_\text{coop} = 0.7$.
The MPPI update finds the minimum-KL perturbation that maximises collective welfare;
otherwise, the consensus is used directly.

\paragraph{Contributions.} This work makes four primary contributions.
\begin{itemize}
    \item We introduce SWA-MPPI, a training-free inference-time alignment framework
    combining WVS-grounded cultural personas, adaptive conflict detection, and
    Prospect-Theory-weighted MPPI optimisation (\S\ref{sec:method}).
    \item We demonstrate mean JSD reductions of 37.6\% (Qwen2.5-72B) and 34.7\%
    (Llama-3.1-70B) versus vanilla baselines across 15 countries
    (\S\ref{sec:results}).
    \item We show that gains are largest for culturally distant country--language
    pairs (Vietnam $+65.8\%$, Mexico $+66.4\%$ for Llama), demonstrating that
    the framework captures genuine cross-cultural signal (\S\ref{sec:results}).
    \item We analyse the adaptive trigger mechanism, showing that selective MPPI
    activation on approximately 35--65\% of scenarios is sufficient to drive
    alignment improvements while maintaining practical latency
    (\S\ref{sec:analysis}).
\end{itemize}

% ============================================================
\section{Related Work}
\label{sec:related}
% ============================================================

\paragraph{LLM moral alignment benchmarks.}
Several benchmarks evaluate whether LLMs reason morally in ways that match human
judgments. \citet{hendrycks2021aligning} propose the ETHICS dataset, covering
commonsense morality, deontology, and utilitarianism, but in English only with
a small crowdsourced ground truth. Commonsense Norm Bank~\citep{commonsense2021}
provides 1.7~million judgments but lacks systematic parametric variation.
MoralExceptQA~\citep{jin2022exceptions} and MoCa~\citep{nie2023moca} introduce
more controlled variation but remain English-centric. PRISM~\citep{kirk2024prism}
and GlobalOpinionQA~\citep{durmus2023global} begin to address cross-cultural
diversity but cover only 4 and handful of languages respectively. MultiTP
\citep{jin2025multitp} is the most comprehensive, supporting 107 languages and
grounding evaluation in 40~million human judgments-we adopt it as our primary
benchmark.

\paragraph{Cross-cultural LLM evaluation.}
A growing body of work examines how LLMs represent different cultural groups.
\citet{santurkar2023opinions} show that LLM outputs are most correlated with the
opinions of liberal, educated, American respondents. \citet{durmus2023global}
observe that alignment with subjective global opinions degrades for non-English
queries. \citet{ryan2024unintended} identify unintended cultural impacts of
RLHF alignment, while \citet{manvi2024geographic} demonstrate geographic bias in
geospatial reasoning. In the moral domain specifically, \citet{takemoto2024moral}
study LLM responses to trolley problems but without cross-linguistic analysis.
Our work is most closely related to \citet{jin2025multitp}, which we extend by
proposing an alignment \emph{intervention} rather than solely a diagnostic study.

\paragraph{Inference-time steering and control.}
Methods for steering LLM outputs at inference time without retraining have gained
traction as an alternative to RLHF. Activation steering~\citep{arditi2025refusal}
manipulates internal representations to modify model behaviour; representation
engineering~\citep{zou2023representation} identifies directions in activation space
corresponding to desired properties. In the prompting space, Constitutional
AI~\citep{bai2022constitutional} uses self-critique loops, and chain-of-thought
prompting~\citep{wei2022chain} elicits structured reasoning. Most relevant to our
work, \citet{sorensen2024roadmap} propose a roadmap for \emph{pluralistic
alignment}, advocating for systems that model the diversity of human values rather
than collapsing to a single preference. SWA-MPPI operationalises this goal through
an ensemble of culturally grounded agents and a stochastic optimal control step.

\paragraph{Model Predictive Path Integral control.}
MPPI~\citep{williams2017mppi} is a sampling-based stochastic optimal control
algorithm that replaces gradient computation with Monte Carlo sampling of
perturbations. It has been applied to robotics~\citep{williams2017mppi} and, more
recently, to token-level generation steering in LLMs~\citep{huang2024mppi}.
Our work is the first to apply MPPI at the logit level for cross-cultural moral
alignment, combining it with Prospect Theory~\citep{kahneman1979prospect} to
model the asymmetric gain-loss sensitivity observed in human moral judgment.

% ============================================================
\section{Method: SWA-MPPI}
\label{sec:method}
% ============================================================

\subsection{Problem Formulation}
\label{sec:method:problem}

Let $\mathcal{D}$ be the MultiTP dataset of moral dilemma scenarios, each of which
requires an LLM to choose between two groups (encoded as decision tokens
\texttt{LEFT}~$\ell$ and \texttt{RIGHT}~$r$). For a target country $c$, let
$\mathbf{h}^{(c)} \in [0,1]^6$ be the human AMCE preference vector encoding the
fraction of human respondents in country $c$ who prefer the normative choice on each
of the six moral dimensions $d \in \{\text{Species, Gender, Age, Fitness,
SocialValue, Utilitarianism}\}$.

A vanilla LLM $f_\theta$ produces logits
$\mathbf{z}(\mathbf{x}) = [z_\ell, z_r] \in \mathbb{R}^2$ for input scenario
$\mathbf{x}$. The probability assigned to sparing the preferred group is
$p_\text{spare} = \sigma(\Delta / T_\text{dec})$, where $\Delta = z_r - z_\ell$
and $T_\text{dec}$ is a decision temperature. Our goal is to find, without
modifying $\theta$, a logit correction $\delta^*$ such that the model AMCE
$\mathbf{m}^{(c)}$ computed from $\sigma((\Delta + \delta^*)/T_\text{dec})$ is
closer to $\mathbf{h}^{(c)}$ under Jensen-Shannon divergence.

\subsection{Cultural Persona Construction from WVS Data}
\label{sec:method:personas}

For each target country $c$, SWA-MPPI constructs $N=4$ persona agents whose
system prompts are grounded in World Values Survey (WVS) Wave~7 statistics.
Three personas correspond to age cohorts (18--35, 36--55, 56+), each parameterised
by WVS variables measuring religiosity, gender egalitarianism, personal autonomy,
meritocratic orientation, and permissiveness toward contextual moral exceptions.
A fourth persona is a utilitarian thinker (always save the greater number of lives).
For countries not covered by WVS, manually authored fallback personas are used.

Concretely, each persona is a natural-language system prompt of the form:

\begin{quote}
\small
\emph{You are a [age-cohort] adult from [country]. Based on the cultural values
of your society, you are [religiosity], [gender egalitarianism], [autonomy], and
[meritocratic]. When making moral judgments you tend to [value description].}
\end{quote}

\noindent The use of WVS data ensures that persona descriptions reflect empirically
measured within-country variation rather than stereotypes or uniform national
identities. Personas are constructed once per country and reused across all
scenarios.

\subsection{Batched Agent Evaluation}
\label{sec:method:eval}

Given scenario $\mathbf{x}$, SWA-MPPI constructs $N+1$ input sequences by
prepending the base system prompt and each of the $N$ persona prefixes to the
query. All $N+1$ sequences are padded to the same length and evaluated in a
\emph{single} batched forward pass, extracting decision logits at the final
token position:
\begin{equation}
    \mathbf{z}_\text{base},\, \mathbf{z}_1, \ldots, \mathbf{z}_N
    = f_\theta(\mathbf{x};\, \varnothing,\, \mathbf{s}_1, \ldots, \mathbf{s}_N)
    \label{eq:batch}
\end{equation}
where each $\mathbf{z}_i = [z_{i,\ell},\, z_{i,r}] / T_\text{cat}$ is temperature-scaled
by a per-category logit temperature $T_\text{cat}$ (see Table~\ref{tab:hyperparams}).

Agent decision gaps are defined as $\delta_i = z_{i,r} - z_{i,\ell}$, and the
consensus signal is their mean:
\begin{equation}
    \bar{\delta} = \frac{1}{N}\sum_{i=1}^{N} \delta_i.
    \label{eq:consensus}
\end{equation}
Agent reward signals relative to the base model are:
\begin{equation}
    r_i = \delta_i - \delta_\text{base}, \qquad i = 1, \ldots, N.
    \label{eq:reward}
\end{equation}

\subsection{Adaptive Conflict Detection}
\label{sec:method:tau}

When agents agree, no optimisation step is needed, and the consensus~$\bar\delta$
is used directly. When agents disagree substantially, the vanilla model's output
may not reflect the target country's balance of views, and an MPPI correction is
warranted. We operationalise this via a per-country adaptive threshold $\tau^{(c)}$
calibrated on a small set of $n_\text{cal} = 50$ scenarios such that the \emph{trigger
rate} (fraction of scenarios activating MPPI) approximates a target rate
$\rho_\text{target} = 0.35$:
\begin{equation}
    \tau^{(c)} = \text{Percentile}_{(1-\rho_\text{target})\times 100}
    \!\left(\left\{\operatorname{Var}(\boldsymbol\delta^{(s)})\right\}_{s=1}^{n_\text{cal}}\right).
    \label{eq:tau_calibration}
\end{equation}
MPPI is triggered on scenario $\mathbf{x}$ if and only if
$\operatorname{Var}(\boldsymbol\delta) \geq \tau^{(c)}$.

\subsection{Prospect-Theory-Weighted MPPI Optimisation}
\label{sec:method:mppi}

When triggered, SWA-MPPI draws $K = 128$ perturbation samples
$\epsilon_k \sim \mathcal{N}(0, \sigma^2)$ with $\sigma = 0.3$ and
evaluates a collective welfare score for each perturbed decision gap
$\tilde\delta_k = \bar\delta + \epsilon_k$.

The utility of perturbation $\epsilon_k$ for agent $i$ combines a \emph{private}
utility based on agent $i$'s own reward and a \emph{social} utility based on
the average reward of the other agents:
\begin{equation}
    U_i(\epsilon_k) = (1 - \lambda_\text{coop})\,v(r_i\,\tilde\delta_k)
    + \lambda_\text{coop}\,v\!\left(\bar{r}_{-i}\,\tilde\delta_k\right),
    \label{eq:utility}
\end{equation}
where $\bar{r}_{-i} = \frac{1}{N-1}\sum_{j \neq i} r_j$, $\lambda_\text{coop} = 0.7$
controls the cooperation weight, and $v(\cdot)$ is the Prospect Theory value
function of \citet{kahneman1979prospect}:
\begin{equation}
    v(x) = \begin{cases}
        x^\alpha & \text{if } x \geq 0, \\
        -\kappa\,|x|^\beta & \text{if } x < 0,
    \end{cases}
    \label{eq:prospect}
\end{equation}
with $\alpha = \beta = 0.88$ (diminishing sensitivity) and $\kappa = 2.25$
(loss aversion). The collective utility is:
\begin{equation}
    U_\text{total}(\epsilon_k) = \frac{1}{N}\sum_{i=1}^{N} U_i(\epsilon_k)
    - \alpha_\text{KL}\,\frac{\epsilon_k^2}{2\sigma^2},
    \label{eq:total_util}
\end{equation}
where the second term is a KL penalty with weight $\alpha_\text{KL} = 0.05$
that prevents excessively large deviations from the prior.

The MPPI update computes importance-weighted perturbation corrections via a
softmax over sample utilities, then returns the optimal correction:
\begin{equation}
    w_k = \frac{\exp(U_\text{total}(\epsilon_k)\,/\,\beta)}
               {\sum_{k'}\exp(U_\text{total}(\epsilon_{k'})\,/\,\beta)},
    \qquad
    \delta^* = \sum_{k=1}^{K} w_k\,\epsilon_k,
    \label{eq:mppi_update}
\end{equation}
where $\beta = T_\text{dec}$ is the MPPI temperature.

\subsection{Positional Debiasing and Final Decision}
\label{sec:method:debias}

Recent work reports that LLMs suffer from recency bias, preferring options
that appear later in a prompt~\citep{liu2024lost}. To cancel this effect, SWA-MPPI
performs two forward passes for each scenario: one with the original ordering
(Group~A in the LEFT lane, Group~B in the RIGHT lane) and one with the ordering
reversed. Decision gaps from the two passes are averaged before applying the MPPI
correction.

The final probability of the preferred choice is:
\begin{equation}
    p_\text{spare} = \sigma\!\left(\frac{\bar\delta + \delta^*}{T_\text{dec}}\right),
    \quad T_\text{dec} = 0.5.
    \label{eq:final}
\end{equation}

The complete SWA-MPPI inference procedure is summarised in Algorithm~\ref{alg:swa}.

\begin{algorithm}[t]
\caption{SWA-MPPI Inference for Country $c$}
\label{alg:swa}
\begin{algorithmic}[1]
\Require Scenario $\mathbf{x}$, frozen LLM $f_\theta$, WVS personas
    $\{\mathbf{s}_i\}_{i=1}^{N}$, calibrated threshold $\tau^{(c)}$
\State \textbf{Batched forward pass:} evaluate $f_\theta$ with base prompt and all
    $N$ persona prefixes in a single batch $\Rightarrow$ obtain logits
    $\mathbf{z}_\text{base}, \mathbf{z}_1, \ldots, \mathbf{z}_N$
    \Comment{Eq.~\ref{eq:batch}}
\State Compute agent decision gaps $\delta_i = z_{i,r} - z_{i,\ell}$,
    consensus $\bar\delta$, and rewards $r_i$ \Comment{Eqs.~\ref{eq:consensus}--\ref{eq:reward}}
\State Compute inter-agent variance $\nu = \operatorname{Var}(\delta_1, \ldots, \delta_N)$
\If{$\nu \geq \tau^{(c)}$} \Comment{Conflict detected}
    \For{$k = 1, \ldots, K$}
        \State Sample $\epsilon_k \sim \mathcal{N}(0, \sigma^2)$
        \State Compute $U_\text{total}(\epsilon_k)$ via Eqs.~\ref{eq:utility}--\ref{eq:total_util}
    \EndFor
    \State Compute importance weights $\{w_k\}$ and correction $\delta^* = \sum_k w_k \epsilon_k$
    \Comment{Eq.~\ref{eq:mppi_update}}
\Else
    \State $\delta^* \leftarrow 0$
\EndIf
\State Repeat steps 1--3 with swapped LEFT/RIGHT ordering; average $\bar\delta$
    \Comment{Positional debiasing}
\State Return $p_\text{spare} = \sigma((\bar\delta + \delta^*) / T_\text{dec})$
    \Comment{Eq.~\ref{eq:final}}
\end{algorithmic}
\end{algorithm}

% ============================================================
\section{Experiments}
\label{sec:experiments}
% ============================================================

\subsection{Dataset and Evaluation Protocol}
\label{sec:exp:data}

We use the MultiTP benchmark~\citep{jin2025multitp}, which provides 460 trolley
problem vignettes per language across 107 languages. Following the benchmark's
evaluation protocol, each vignette asks an LLM to choose which of two groups a
self-driving car with brake failure should spare. Groups vary across six moral
dimensions (see Table~\ref{tab:categories}).

For each of the 15 target countries, we load the vignette set in the country's
\emph{native language} (see Table~\ref{tab:results_qwen}), apply the deduplication
step (keep \texttt{paraphrase=0} vignettes only), filter low-quality Utilitarianism
scenarios (equal-composition groups), and balance the resulting pool to at most 80
vignettes per category. After augmentation of under-represented categories (Species,
Utilitarianism) to a minimum of 36, each country is evaluated on \textbf{342
scenarios}. Side balance is enforced at 50\% preferred-on-right.

Ground truth is the country-specific AMCE data from the MultiTP repository
(\texttt{country\_specific\_ACME.csv}), which aggregates over 40~million human
responses from \citet{awad2018moral} into six preference scores per country.

\begin{table}[t]
\caption{The six moral dimensions in MultiTP with their default ``preferred'' choice
    and per-category logit temperature used in SWA-MPPI.}
\label{tab:categories}
\centering
\begin{tabular}{llll}
\toprule
Dimension       & Preferred choice  & Contra choice    & Logit $T$ \\
\midrule
Species         & Human             & Animal (dog/cat) & 4.0       \\
Gender          & Female            & Male             & 3.5       \\
Age             & Young             & Elderly          & 1.5       \\
Fitness         & Fit (athlete)     & Unfit (large)    & 1.5       \\
Social Value    & High (executive)  & Low (homeless)   & 1.5       \\
Utilitarianism  & More lives        & Fewer lives      & 1.5       \\
\bottomrule
\end{tabular}
\end{table}

\subsection{Models}
\label{sec:exp:models}

We evaluate on two frontier 70B-class instruction-following models loaded in
4-bit quantisation via Unsloth~\citep{unsloth2024} on a single NVIDIA H100 80\,GB GPU:
\begin{itemize}
    \item \textbf{Meta-Llama-3.1-70B-Instruct} (FP8, \texttt{neuralmagic} quantisation)
    \item \textbf{Qwen2.5-72B-Instruct} (GPTQ Int8, \texttt{unsloth} quantisation)
\end{itemize}
All generation uses greedy decoding (temperature~0) and a fixed random seed (42).
Logits are extracted only for tokens \texttt{LEFT} (id~22397) and \texttt{RIGHT}
(id~11572) at the final position.

\subsection{Baselines}
\label{sec:exp:baseline}

The \textbf{Vanilla LLM} baseline processes each scenario with a single forward
pass using a generic system prompt (\emph{``You are a helpful assistant.''}),
extracts LEFT/RIGHT logits, and computes $p_\text{spare}$ with the same decision
temperature $T_\text{dec} = 0.5$. No persona is active and no MPPI step is taken.
The vanilla baseline is run per-country in the native language of that country,
ensuring a fair comparison.

\subsection{Evaluation Metrics}
\label{sec:exp:metrics}

For each country $c$ and scenario set we compute country-level AMCE vectors from
model predictions and compare them to human AMCE vectors using:
\begin{itemize}
    \item \textbf{JSD}: Jensen-Shannon Distance $\in [0,1]$ between normalised AMCE vectors (primary metric, lower is better).
    \item \textbf{Pearson $r$}: Pearson correlation between model and human AMCE vectors across the 6 dimensions.
    \item \textbf{Cosine similarity}: Cosine similarity between the two AMCE vectors.
    \item \textbf{MAE}: Mean absolute error in percentage points between model and human AMCE values.
    \item \textbf{RMSE}: Root mean square error in percentage points.
    \item \textbf{MPPI trigger rate}: Fraction of scenarios that activated the MPPI step.
    \item \textbf{MPPI flip rate}: Among triggered scenarios, fraction where MPPI reversed the vanilla decision.
\end{itemize}
All metrics are reported per-country and aggregated as unweighted means across the
15 target countries.

\subsection{Hyperparameter Configuration}
\label{sec:exp:hyperparams}

Table~\ref{tab:hyperparams} lists all SWA-MPPI hyperparameters. Values were
not tuned on MultiTP; they were set based on the Prospect Theory literature
($\alpha,\beta,\kappa$) and a preliminary sanity-check on a held-out synthetic
dataset.

\begin{table}[t]
\caption{SWA-MPPI hyperparameter configuration.}
\label{tab:hyperparams}
\centering
\begin{tabular}{lll}
\toprule
Parameter                  & Symbol                  & Value       \\
\midrule
Number of persona agents   & $N$                     & 4           \\
MPPI sample size           & $K$                     & 128         \\
Perturbation std.\ dev.\   & $\sigma$                & 0.3         \\
Cooperation weight         & $\lambda_\text{coop}$   & 0.7         \\
KL penalty weight          & $\alpha_\text{KL}$      & 0.05        \\
PT gain curvature          & $\alpha$                & 0.88        \\
PT loss curvature          & $\beta$                 & 0.88        \\
PT loss aversion           & $\kappa$                & 2.25        \\
Decision temperature       & $T_\text{dec}$          & 0.5         \\
Target trigger rate        & $\rho_\text{target}$    & 0.35        \\
Calibration set size       & $n_\text{cal}$          & 50          \\
\bottomrule
\end{tabular}
\end{table}

% ============================================================
\section{Results}
\label{sec:results}
% ============================================================

\subsection{Overall Alignment Improvement}
\label{sec:results:overall}

Tables~\ref{tab:results_llama} and \ref{tab:results_qwen} report per-country
alignment metrics for Meta-Llama-3.1-70B-Instruct and Qwen2.5-72B-Instruct
respectively. Both tables include vanilla and SWA-MPPI values for JSD, Pearson~$r$,
MAE, and RMSE, alongside the percentage improvement in JSD.

\paragraph{Llama-3.1-70B.}
SWA-MPPI reduces mean JSD from 0.0937 to 0.0612 (\textbf{$-$34.7\%}) and mean
Pearson~$r$ improves from $-0.026$ to $+0.384$ ($+0.41$). Mean MAE drops by 4.2
percentage points (from 17.8 to 13.6\,pp). Improvements are consistent across
most countries, with the largest gains observed for Vietnam ($-65.8\%$ JSD),
Mexico ($-66.4\%$), and Nigeria ($-37.5\%$). The only countries where SWA-MPPI
does not improve are Brazil ($+23.3\%$ JSD increase), where the model already
achieves a relatively low vanilla JSD, and Saudi Arabia, where cultural distances
are large and persona grounding is weaker due to limited WVS coverage.

\paragraph{Qwen2.5-72B.}
SWA-MPPI reduces mean JSD from 0.0839 to 0.0536 (\textbf{$-$37.6\%}) and mean
Pearson~$r$ improves from $+0.348$ to $+0.696$ ($+0.148$). Mean MAE decreases
from 13.7 to 11.0\,pp. Qwen2.5 generally achieves lower vanilla baseline JSD than
Llama (0.084 vs.\ 0.094), yet SWA-MPPI still delivers substantial improvements,
particularly for the USA ($-71.1\%$), Australia ($-75.4\%$), Japan ($-43.9\%$),
and Vietnam ($-39.9\%$). Qwen2.5 appears to benefit more on English-speaking
Anglosphere countries (USA, AUS, GBR) than Llama does, possibly reflecting
differences in the models' underlying training distributions.

\begin{table}[t]
\caption{Per-country alignment results for \textbf{Meta-Llama-3.1-70B-Instruct}.
    Van.\ = Vanilla baseline; SWA = SWA-MPPI. $\uparrow$/$\downarrow$ indicate
    the direction of improvement. JSD\% = percentage improvement in JSD
    (positive = improved).}
\label{tab:results_llama}
\centering
\setlength{\tabcolsep}{4pt}
\footnotesize
\begin{tabular}{lcccccccccc}
\toprule
\multirow{2}{*}{Country} & \multirow{2}{*}{Lang.}
  & \multicolumn{2}{c}{JSD $\downarrow$}
  & \multicolumn{2}{c}{Pearson $r$ $\uparrow$}
  & \multicolumn{2}{c}{MAE $\downarrow$}
  & \multicolumn{2}{c}{RMSE $\downarrow$}
  & JSD\% \\
\cmidrule(lr){3-4}\cmidrule(lr){5-6}\cmidrule(lr){7-8}\cmidrule(lr){9-10}
 & & Van. & SWA & Van. & SWA & Van. & SWA & Van. & SWA & \\
\midrule
USA & en    & 0.156 & 0.123 & 0.012 & 0.197 & 19.5 & 15.3 & 26.2 & 21.8 & +20.9\% \\
DEU & de    & 0.071 & 0.039 & $-$0.455 & 0.685 & 14.4 & 12.5 & 17.8 & 14.5 & +44.9\% \\
CHN & zh    & 0.046 & 0.031 & 0.406 & 0.802 & 16.7 & 12.7 & 18.8 & 14.0 & +32.3\% \\
JPN & ja    & 0.046 & 0.042 & 0.496 & 0.466 & 15.9 & ~9.8 & 18.0 & 12.4 & $+$9.0\% \\
BRA & pt    & 0.053 & 0.066 & 0.349 & 0.348 & 12.8 & 10.4 & 15.8 & 12.8 & $-$23.3\% \\
SAU & ar    & 0.054 & 0.060 & 0.043 & $-$0.210 & 15.8 & 17.2 & 18.7 & 20.2 & $-$10.7\% \\
VNM & vi    & 0.155 & 0.053 & 0.519 & 0.136 & 21.2 & 14.2 & 26.1 & 16.8 & +65.8\% \\
FRA & fr    & 0.050 & 0.052 & 0.617 & 0.394 & 17.3 & 17.0 & 19.6 & 19.4 & $-$3.8\% \\
IND & hi    & 0.044 & 0.034 & 0.630 & 0.715 & 16.7 & 14.3 & 18.6 & 15.7 & +23.6\% \\
KOR & ko    & 0.089 & 0.044 & $-$0.272 & 0.551 & 19.1 & 12.0 & 23.7 & 14.7 & +51.2\% \\
GBR & en    & 0.157 & 0.092 & 0.012 & 0.286 & 19.7 & 13.8 & 26.5 & 17.8 & +41.6\% \\
RUS & ru    & 0.045 & 0.041 & 0.539 & 0.637 & 15.7 & 12.6 & 17.8 & 14.5 & +17.9\% \\
MEX & es    & 0.127 & 0.043 & $-$0.371 & 0.307 & 24.0 & 15.2 & 28.9 & 16.9 & +66.4\% \\
NGA & en    & 0.158 & 0.099 & 0.109 & 0.497 & 19.7 & 13.8 & 26.7 & 20.0 & +37.5\% \\
AUS & en    & 0.155 & 0.101 & 0.039 & 0.224 & 19.4 & 13.9 & 26.1 & 18.9 & +34.9\% \\
\midrule
\textbf{Mean} & & \textbf{0.094} & \textbf{0.061} & $-$0.026 & 0.384 & 17.8 & 13.6 & 21.9 & 16.7 & \textbf{+34.7\%} \\
\bottomrule
\end{tabular}
\end{table}

\begin{table}[t]
\caption{Per-country alignment results for \textbf{Qwen2.5-72B-Instruct}.
    Notation as in Table~\ref{tab:results_llama}.}
\label{tab:results_qwen}
\centering
\setlength{\tabcolsep}{4pt}
\footnotesize
\begin{tabular}{lcccccccccc}
\toprule
\multirow{2}{*}{Country} & \multirow{2}{*}{Lang.}
  & \multicolumn{2}{c}{JSD $\downarrow$}
  & \multicolumn{2}{c}{Pearson $r$ $\uparrow$}
  & \multicolumn{2}{c}{MAE $\downarrow$}
  & \multicolumn{2}{c}{RMSE $\downarrow$}
  & JSD\% \\
\cmidrule(lr){3-4}\cmidrule(lr){5-6}\cmidrule(lr){7-8}\cmidrule(lr){9-10}
 & & Van. & SWA & Van. & SWA & Van. & SWA & Van. & SWA & \\
\midrule
USA & en    & 0.107 & 0.031 & 0.400 & 0.822 & 16.3 & ~9.3 & 22.7 & 10.8 & +71.1\% \\
DEU & de    & 0.060 & 0.060 & 0.845 & 0.747 & ~9.7 & ~9.6 & 11.8 & 12.0 & ~+0.0\% \\
CHN & zh    & 0.051 & 0.037 & 0.851 & 0.762 & ~9.2 & ~9.1 & ~9.6 & 10.9 & +27.8\% \\
JPN & ja    & 0.099 & 0.055 & 0.634 & 0.786 & 15.6 & 13.6 & 19.5 & 15.7 & +43.9\% \\
BRA & pt    & 0.075 & 0.070 & 0.349 & 0.529 & 12.6 & 12.1 & 15.9 & 14.5 & $-$6.7\% \\
SAU & ar    & 0.078 & 0.080 & 0.443 & 0.527 & 11.5 & 13.5 & 14.2 & 14.7 & $-$2.9\% \\
VNM & vi    & 0.130 & 0.078 & 0.559 & 0.456 & 16.1 & 11.8 & 21.8 & 15.3 & +39.9\% \\
FRA & fr    & 0.079 & 0.048 & 0.617 & 0.838 & 14.3 & ~9.0 & 16.2 & 10.7 & +39.4\% \\
IND & hi    & 0.060 & 0.056 & 0.670 & 0.739 & 11.0 & 11.2 & 12.2 & 13.5 & $+$6.9\% \\
KOR & ko    & 0.072 & 0.051 & 0.355 & 0.645 & 15.3 & 12.8 & 18.2 & 14.0 & +28.8\% \\
GBR & en    & 0.106 & 0.029 & 0.419 & 0.800 & 16.5 & ~8.2 & 22.7 & ~9.9 & +72.2\% \\
RUS & ru    & 0.080 & 0.063 & 0.623 & 0.784 & 12.4 & 10.1 & 15.1 & 12.6 & +21.4\% \\
MEX & es    & 0.076 & 0.053 & 0.606 & 0.545 & 12.5 & 10.4 & 14.7 & 11.0 & +29.8\% \\
NGA & en    & 0.110 & 0.065 & 0.407 & 0.615 & 16.5 & 15.4 & 23.3 & 18.8 & +41.2\% \\
AUS & en    & 0.105 & 0.026 & 0.438 & 0.841 & 16.4 & ~9.0 & 22.4 & 10.2 & +75.4\% \\
\midrule
\textbf{Mean} & & \textbf{0.084} & \textbf{0.054} & 0.348 & 0.696 & 13.7 & 11.0 & 17.3 & 13.0 & \textbf{+37.6\%} \\
\bottomrule
\end{tabular}
\end{table}

\subsection{Per-Dimension AMCE Analysis}
\label{sec:results:amce}

To understand \emph{which} moral dimensions benefit most from SWA-MPPI, we compute
the mean absolute bias per dimension, averaged across countries:
\begin{equation}
    \bar{e}_d = \frac{1}{|C|}\sum_{c \in C} |m_d^{(c)} - h_d^{(c)}|,
    \label{eq:dimension_mae}
\end{equation}
where $m_d^{(c)}$ and $h_d^{(c)}$ are model and human AMCE values for dimension $d$
in country~$c$.

For Qwen2.5-72B, the largest improvements (vanilla $\to$ SWA) are on
\textbf{SocialValue} ($-8.2$\,pp), \textbf{Utilitarianism} ($-6.4$\,pp), and
\textbf{Age} ($-5.1$\,pp). These are the dimensions where the vanilla model tends
to produce overconfident, binarised outputs (e.g., near-100\% preference for high
social status) that diverge most from the graded human distributions. In contrast,
\textbf{Species} shows the smallest gain, since both vanilla models and humans
already strongly prefer sparing humans over animals, leaving little room for
improvement.

\subsection{Country-Level Analysis}
\label{sec:results:country}

Both models exhibit consistent patterns across countries. Countries in the
\emph{Global East} (CHN, JPN, KOR) generally achieve moderate vanilla alignment
and benefit moderately from SWA-MPPI, consistent with \citet{jin2025multitp}'s
finding that East Asian AMCE profiles are moderately distinct from the Western
values encoded in model training data. Countries in the \emph{Global South} and
those with limited WVS representation (SAU, BRA) show less consistent improvement,
highlighting a key limitation: persona quality degrades when WVS data are sparse
or absent (see \S\ref{sec:analysis:limitations}).

English-speaking countries (USA, GBR, AUS) show some of the largest vanilla JSD
values despite the models being trained on predominantly English data. This is
counter-intuitive but consistent with \citet{jin2025multitp}: RLHF alignment
introduces specific moral biases (e.g., extreme protection of animals, rigid
gender neutrality) that deviate from average Anglosphere human preferences.
SWA-MPPI substantially corrects these biases ($-71\%$ to $-75\%$ JSD for Qwen2.5
on USA, GBR, and AUS), suggesting that the WVS-derived utilitarian and age-cohort
personas successfully counterbalance RLHF-induced overconfidence.

% ============================================================
\section{Analysis}
\label{sec:analysis}
% ============================================================

\subsection{Adaptive MPPI Trigger Mechanism}
\label{sec:analysis:trigger}

Table~\ref{tab:trigger} reports MPPI trigger rates and flip rates per country for
Qwen2.5-72B. Trigger rates range from 47\% (AUS) to 65\% (JPN), reflecting genuine
inter-agent disagreement variability across countries and categories. Critically,
\emph{flip rates}-the fraction of triggered scenarios where MPPI reverses the
vanilla decision-are low (1.6\%--14.3\%), indicating that MPPI primarily
\emph{modulates} decision confidence rather than reversing decisions wholesale.
This is desirable: strong vanilla decisions are preserved while ambiguous ones are
resolved in a culturally informed direction.

\begin{table}[t]
\caption{MPPI trigger and flip rates per country for Qwen2.5-72B-Instruct.
    Flip rate is the fraction of \emph{triggered} scenarios where MPPI reversed
    the vanilla decision direction.}
\label{tab:trigger}
\centering
\setlength{\tabcolsep}{5pt}
\footnotesize
\begin{tabular}{lcccc}
\toprule
Country & Trigger rate & Flip rate & Mean variance & Calibrated $\tau$ \\
\midrule
USA & 53.5\% & 4.9\% & 1.462 & 0.512 \\
DEU & 62.3\% & 4.2\% & 1.437 & 0.156 \\
CHN & 57.0\% & 10.8\% & 1.803 & 0.470 \\
JPN & 64.9\% & 7.2\% & 4.257 & 1.388 \\
BRA & 57.6\% & 6.6\% & 2.416 & 0.604 \\
SAU & 59.4\% & 14.3\% & 1.064 & 0.381 \\
VNM & 56.4\% & 1.6\% & 1.152 & 0.164 \\
FRA & 54.7\% & 3.8\% & 1.501 & 0.487 \\
IND & 50.9\% & 5.8\% & 1.298 & 0.442 \\
KOR & 55.3\% & 6.7\% & 1.632 & 0.512 \\
GBR & 55.6\% & 3.7\% & 1.721 & 0.575 \\
RUS & 49.4\% & 4.3\% & 1.448 & 0.512 \\
MEX & 58.2\% & 5.1\% & 1.588 & 0.438 \\
NGA & 56.1\% & 8.4\% & 1.623 & 0.481 \\
AUS & 47.1\% & 3.1\% & 1.644 & 0.792 \\
\midrule
\textbf{Mean} & 56.0\% & 6.0\% & 1.769 & 0.494 \\
\bottomrule
\end{tabular}
\end{table}

\subsection{Latency and Computational Cost}
\label{sec:analysis:latency}

Mean per-scenario latency for SWA-MPPI is 730\,ms on Qwen2.5-72B (H100 80\,GB),
compared to 160\,ms for the vanilla single-pass baseline. The $4.6\times$ overhead
is dominated by two factors: (i) the batched multi-agent forward pass ($N+1 = 5$
inputs vs.\ 1), and (ii) the debiasing second pass. The MPPI sample loop
($K = 128$) operates on cached logit tensors and contributes negligibly to wall-clock
time. At 342 scenarios per country and 15 countries, total evaluation time per model
is approximately 62 minutes.

\subsection{Sensitivity to Cooperation Weight $\lambda_\text{coop}$}
\label{sec:analysis:ablation_lambda}

The cooperation weight $\lambda_\text{coop}$ governs the balance between each
agent's private cultural utility and the social welfare of the group. The
configuration range spans $\lambda_\text{coop} \in \{0.0, 0.2, 0.4, 0.6, 0.8, 1.0\}$.
At $\lambda = 0.0$ (pure private utility), agents act selfishly and the MPPI
optimisation degenerates to maximising the strongest persona's preference, which
tends to over-correct for outlier personas. At $\lambda = 1.0$ (pure social utility),
the method collapses to the equally-weighted consensus and MPPI optimisation
becomes trivial. The value $\lambda_\text{coop} = 0.7$ was found to provide the
best trade-off, giving sufficient weight to social cohesion while preserving
persona-specific sensitivity.

\subsection{Limitations}
\label{sec:analysis:limitations}

\paragraph{WVS coverage.}
13 of 15 target countries have WVS Wave~7 coverage; Saudi Arabia and Vietnam rely
on partial coverage. Countries outside the WVS (e.g., many African nations) would
require manually authored personas, reducing grounding quality.

\paragraph{Binary decision format.}
The MultiTP benchmark presents strictly binary choices. SWA-MPPI is designed for
this setting; extending the framework to multi-option or continuous preference
elicitation is left for future work.

\paragraph{Language-culture mapping.}
Following \citet{jin2025multitp}, we use one language per country. Countries with
multiple dominant languages (e.g., India, Nigeria) are represented by a single
language variety, potentially missing intra-national cultural diversity.

\paragraph{Model quantisation.}
Both models are loaded in 4-bit or Int8 quantisation. While quantisation loss is
typically small for 70B models, it may introduce systematic logit artefacts that
affect the MPPI optimisation.

% ============================================================
\section{Conclusion}
\label{sec:conclusion}
% ============================================================

We have presented SWA-MPPI, a training-free inference-time framework for
cross-cultural moral alignment of large language models. The method constructs
culturally grounded persona agents from World Values Survey data, evaluates their
logit-level consensus in a single batched forward pass, and applies a
Prospect-Theory-weighted Model Predictive Path Integral correction when
inter-agent disagreement signals that the model's vanilla output may not reflect
the target population's moral balance.

Evaluated on the MultiTP benchmark across 15 countries, 12 native languages, and
6 moral dimensions, SWA-MPPI reduces Jensen-Shannon Distance to human moral
preferences by 37.6\% on Qwen2.5-72B and 34.7\% on Llama-3.1-70B, while improving
Pearson correlation with human AMCE vectors by up to $+0.41$ points. Improvements
are largest for culturally distant country--language pairs, confirming that the
framework captures genuine cross-cultural signal rather than artefacts of language
surface form.

SWA-MPPI operates entirely on frozen model weights, requires no labelled preference
data for training, and introduces a mean per-scenario overhead of approximately
730\,ms on a single H100 GPU-a practical cost for deployment-scale culturally
aware inference. We hope this work contributes a concrete step toward the pluralistic
alignment vision of \citet{sorensen2024roadmap}: AI systems that genuinely model
the diversity of human moral preferences rather than imposing a single culturally
unanchored stance.

\paragraph{Future work.}
Promising directions include extending the persona library to cover a larger set
of WVS variables and more languages, investigating the interaction of SWA-MPPI
with fine-tuned models, integrating the framework with structured chain-of-thought
prompting for transparent moral reasoning, and applying the MPPI paradigm to
multi-option preference scenarios beyond trolley problems.

% ============================================================
\begin{ack}
The authors acknowledge the creators of the MultiTP dataset~\citep{jin2025multitp}
and the Moral Machine experiment~\citep{awad2018moral} for making their data
publicly available. Experiments were conducted on NVIDIA H100 hardware.
\end{ack}

% ============================================================
\bibliographystyle{abbrvnat}
\bibliography{references}

\end{document}